\scanpage{3}
\mspage{( 2 )}

Soient $\mathrm A$ un anneau local nœthérien et $\widehat{\mathrm A}$ son complété. On a :

\textsc{Lemme.} --- \textit{Le noyau de l’application canonique : $\widehat{\mathrm A}\underset{\mathrm A}{\otimes}\mathrm D(\mathrm A)\to\mathrm D(\widehat{\mathrm A})$ est contenu dans l’intersection : $\displaystyle\bigcap_{\mathrm M\in\sigma}\left(\widehat{\mathrm A}\underset{\mathrm A}{\otimes}\mathrm M\right)$, où $\sigma$ désigne l’ensemble des sous-modules $\mathrm M$ de $\mathrm D(\mathrm A)$ tels que $\mathrm F=\mathrm D(\mathrm A)/\mathrm M$ soit un $\mathrm A$-module de type fini.}

De ce lemme on déduit que si $\mathrm D(\mathrm A)$ est semi-fini au sens de Suzuki (\textsuperscript{10}), c’est-à-dire somme directe d’un module libre et d’un module de type fini, la suite canonique
\[
0\to\widehat{\mathrm A}\underset{\mathrm A}{\otimes}\mathrm D(\mathrm A)\to(\widehat{\mathrm A})\to\mathrm D_{\mathrm A}(\widehat{\mathrm A})\to0
\]
est exacte.

Cela entraîne notamment :

\textsc{Proposition 4.} --- \textit{Si $\mathrm A$ est une localité analytique\uline{ment intègre}, $\mathrm K_{\widehat{\mathrm A}}$ est séparable sur $\mathrm K_{\mathrm A}$.}

De même (\textit{cf.} cor. 1 de la prop. 1), si $\mathrm A$ est un anneau de séries convergentes sur un corps valué complet de caractéristique $p>0$ et tel que $(k:k^p)<+\infty$, $\mathrm K_{\widehat{\mathrm A}}$ est séparable sur $\mathrm K_{\mathrm A}$. Ceci permet de combler ce qui nous paraît être une lacune dans la démonstration donnée par Lang (\textsuperscript{4}) du résultat suivant :
% slanted
\margin{\handwritten{$\left.\rule{0pt}{8em}\right]$ connu}}

\textsc{Proposition 5.} --- \textit{Si $\mathrm A=k\{\mathrm T\}$ est l’anneau des séries convergentes sur un corps $k$ valué complet et algébriquement clos, $\mathrm K_{\mathrm A}$ est quasialgébriquement clos $(\mathrm C^1)$.}

Un cas particulier de la proposition 4, classique, est que $\mathrm K((\mathrm X_1,\mathrm X_2,\ldots,\mathrm X_n))$ est séparable sur $\mathrm K(\mathrm X_1,\ldots,\mathrm X_n)$ pour tout corps $\mathrm K$. Il implique, en tenant compte du théorème de structure de Cohen :

\textsc{Proposition 6.} --- \textit{Soit $\mathrm A$ un anneau local régulier contenant un corps $k$ tel que $\mathrm L_{\mathrm A}$ soit séparable sur $k$. Alors $\mathrm K_{\mathrm A}$ est aussi séparable sur $k$.}

3. \textsc{Critère de réduction analytique.} --- La comparaison de la proposition 1 et des énoncés de Suzuki et Mount suggère l’existence d’un critère de réduction analytique sous-jacent; nous nous proposons de voir dans ce paragraphe comment les conclusions du paragraphe 2 permettent d’en obtenir un. Notons tout d’abord le lemme suivant qui peut être considéré comme classique :

\textsc{Lemme.} --- \textit{Soit $\mathrm A$ un anneau local nœthérien analytiquement intègre et tel que $\mathrm K_{\widehat{\mathrm A}}$ soit séparable sur $\mathrm K_{\mathrm A}$. Si $\mathrm B$ est un anneau local intègre de la forme $\mathrm B=\mathrm C_{\mathfrak p}$, où $\mathrm C$ est une $\mathrm A$-algèbre \uline{finie} et $\mathfrak p$ un idéal premier de $\mathrm C$ tel que $\mathfrak p\cap\mathrm A=\mathfrak M$ ($\mathfrak M$ idéal maximal de $\mathrm A$), $\mathrm B$ est analytiquement réduit.}
% slanted
\margin{\handwritten{$\big|$ faux si on ne \unsure{suppose} pas $A\to B$ injectif ??}}

On en déduit :

\textsc{Théorème.} --- \textit{Soit $\mathrm A$ un anneau local régulier vérifiant l’une ou l’autre des conditions suivantes :}

$(h_1)$ \textit{$\mathrm K_{\mathrm A}$ est de caractéristique $0$;}

$(h_2)$ \tikz[baseline=(x.base)]{\node[draw,ellipse,inner sep=2pt] (x) {\itshape $\mathrm D(\mathrm A)$ est libre.};}
% archive sheet number: 2
